\documentclass[10pt,a4paper]{amsart}
\usepackage{amsmath,amssymb,amsthm,graphicx}
% Retained minimal preamble for the Stacks Project chapter more-algebra.tex (More on Algebra).
% The chapter file pinned-original-source.tex carries only \input{preamble} on its first line, so no
% preamble text lives inside the chapter itself.  The Stacks Project uses one single shared preamble for
% every chapter; the snapshot stages it identically as sources/stacks-chow/preamble.tex and
% sources/stacks-etale-cohomology/preamble.tex.  The block below is copied mechanically, line for line,
% from sources/stacks-chow/preamble.tex lines 14-18 (the Xy-pic setup) and 154-217 (the native theorem
% declarations together with the mathematical macros).  Nothing was added, removed or redefined; the
% wrapper class and the amsmath/amssymb/amsthm/graphicx base packages are supplied by the assembler.
% Chapter: more-algebra.tex -- More on Algebra.
% For commutative diagrams we use Xy-pic
\usepackage[all]{xy}

% We use 2cell for 2-commutative diagrams.
\xyoption{2cell}

% Theorem environments.
%
\theoremstyle{plain}
\newtheorem{theorem}[subsection]{Theorem}
\newtheorem{proposition}[subsection]{Proposition}
\newtheorem{lemma}[subsection]{Lemma}

\theoremstyle{definition}
\newtheorem{definition}[subsection]{Definition}
\newtheorem{example}[subsection]{Example}
\newtheorem{exercise}[subsection]{Exercise}
\newtheorem{situation}[subsection]{Situation}

\theoremstyle{remark}
\newtheorem{remark}[subsection]{Remark}
\newtheorem{remarks}[subsection]{Remarks}

\numberwithin{equation}{subsection}

% Macros
%
\def\lim{\mathop{\mathrm{lim}}\nolimits}
\def\colim{\mathop{\mathrm{colim}}\nolimits}
\def\Spec{\mathop{\mathrm{Spec}}}
\def\Hom{\mathop{\mathrm{Hom}}\nolimits}
\def\Ext{\mathop{\mathrm{Ext}}\nolimits}
\def\SheafHom{\mathop{\mathcal{H}\!\mathit{om}}\nolimits}
\def\SheafExt{\mathop{\mathcal{E}\!\mathit{xt}}\nolimits}
\def\Sch{\mathit{Sch}}
\def\Mor{\mathop{\mathrm{Mor}}\nolimits}
\def\Ob{\mathop{\mathrm{Ob}}\nolimits}
\def\Sh{\mathop{\mathit{Sh}}\nolimits}
\def\NL{\mathop{N\!L}\nolimits}
\def\CH{\mathop{\mathrm{CH}}\nolimits}
\def\proetale{{pro\text{-}\acute{e}tale}}
\def\etale{{\acute{e}tale}}
\def\QCoh{\mathit{QCoh}}
\def\Ker{\mathop{\mathrm{Ker}}}
\def\Im{\mathop{\mathrm{Im}}}
\def\Coker{\mathop{\mathrm{Coker}}}
\def\Coim{\mathop{\mathrm{Coim}}}

% Boxtimes
%
\DeclareMathSymbol{\boxtimes}{\mathbin}{AMSa}{"02}

%
% Macros for moduli stacks/spaces
%
\def\QCohstack{\mathcal{QC}\!\mathit{oh}}
\def\Cohstack{\mathcal{C}\!\mathit{oh}}
\def\Spacesstack{\mathcal{S}\!\mathit{paces}}
\def\Quotfunctor{\mathrm{Quot}}
\def\Hilbfunctor{\mathrm{Hilb}}
\def\Curvesstack{\mathcal{C}\!\mathit{urves}}
\def\Polarizedstack{\mathcal{P}\!\mathit{olarized}}
\def\Complexesstack{\mathcal{C}\!\mathit{omplexes}}
% \Pic is the operator that assigns to X its picard group, usage \Pic(X)
% \Picardstack_{X/B} denotes the Picard stack of X over B
% \Picardfunctor_{X/B} denotes the Picard functor of X over B
\def\Pic{\mathop{\mathrm{Pic}}\nolimits}
\def\Picardstack{\mathcal{P}\!\mathit{ic}}
\def\Picardfunctor{\mathrm{Pic}}
\def\Deformationcategory{\mathcal{D}\!\mathit{ef}}

% --- appended 2026-10-12 (Chief): the project own definitions of the reference/slogan/history
% display environments, copied mechanically from the shared Stacks preamble lines 8-12
% (\usepackage{verbatim} plus the three \newenvironment lines, which the chapter preamble keeps
% verbatim from the project preamble). Nothing added, removed or redefined: these environments
% carry no mathematical content and the project itself typesets them as comments. Without these
% three lines a counted statement whose own text uses them fails with Environment slogan undefined.
\usepackage{verbatim}
\newenvironment{reference}{\comment}{\endcomment}
\newenvironment{slogan}{\comment}{\endcomment}
\newenvironment{history}{\comment}{\endcomment}
\usepackage{cleveref}
\crefname{originaldefinition}{original definition}{original definitions}
\Crefname{originaldefinition}{Original definition}{Original definitions}
\makeatletter
\newlabel{derived-definition-derived-colimit}{{090Z}{source}{}{sourcepointer-1}{}}
\newlabel{derived-definition-derived-colimit@cref}{{[originaldefinition][0][]090Z}{[1][0][]source}{}{}{}}
\makeatother
\crefname{originallemma}{original lemma}{original lemmas}
\Crefname{originallemma}{Original lemma}{Original lemmas}
\makeatletter
\newlabel{lemma-dual-perfect-complex}{{07VI}{source}{}{sourcepointer-2}{}}
\newlabel{lemma-dual-perfect-complex@cref}{{[originallemma][0][]07VI}{[1][0][]source}{}{}{}}
\makeatother
\crefname{originalequation}{original equation}{original equations}
\Crefname{originalequation}{Original equation}{Original equations}
\makeatletter
\newlabel{equation-internal-hom}{{0A63}{source}{}{sourcepointer-3}{}}
\newlabel{equation-internal-hom@cref}{{[originalequation][0][]0A63}{[1][0][]source}{}{}{}}
\makeatother
\title{Trivial duality for systems of perfect objects: RHom(hocolim K\_n, E) = Rlim E tensor K\_n\textasciicircum{}dual}
\author{The Stacks Project Authors}
\date{stacks-project:more-algebra.tex:local-snapshot-20261012; GNU FDL 1.2 or later}
\begin{document}
\maketitle
\setcounter{subsection}{1}
\noindent\textit{Editorial note.} All delivered mathematical text is the author's own, copied from the pinned Stacks Project chapter file more-algebra.tex; no statement or proof step was rewritten, reordered or removed. The counted span is the author environment and its own proof environment. References inside the span are handed to the assembler as navigation stubs that carry the canonical tag of the cited tag only; whether each of them is load-bearing is recorded on the stub itself.\par\medskip
\noindent\textit{Allowed standard prerequisites (editorial).}\begin{itemize}
\item For a perfect complex K\_n with dual K\_n\textasciicircum{}dual = RHom(K\_n, A) one has E tensor K\_n\textasciicircum{}dual = RHom\_A(K\_n, E) for every E. invoked at proof lines 19251-19253. Turns the right-hand side into a limit of Hom complexes.
\item The derived colimit of a system is the cone of the canonical map from the direct sum to itself. invoked at statement line 19242. Identifies K with the cone of a map of direct sums, matching the Rlim triangle.
\item The internal-hom adjunction formula Hom(K, RHom(L, M)) = Hom(K tensor L, M). invoked at proof lines 19263-19264. Reduces Hom out of a direct sum to a product of Hom complexes.
\end{itemize}\medskip
\ifdefined\hypertarget\hypertarget{sourcepointer-1}{}\fi
\noindent\textit{Source reference.} Original definition 090Z: sources/stacks-chow/tags.txt, source lines 11427--11427. Cited at statement line 19242 to define the derived colimit K = hocolim K\_n; the definition is used inside the counted span.\par\medskip
\ifdefined\hypertarget\hypertarget{sourcepointer-2}{}\fi
\noindent\textit{Source reference.} Original lemma 07VI: sources/stacks-more-algebra/pinned-original-source.tex, source lines 19176--19176. Cited at proof line 19251 to rewrite E tensor K\_n\textasciicircum{}dual as RHom\_A(K\_n, E); the conclusion is used inside the counted span.\par\medskip
\ifdefined\hypertarget\hypertarget{sourcepointer-3}{}\fi
\noindent\textit{Source reference.} Original equation 0A63: sources/stacks-more-algebra/pinned-original-source.tex, source lines 18654--18654. Cited at proof line 19263 to reduce the Hom out of the direct sum to a product of Hom complexes; the formula is used inside the counted span.\par\medskip
\begin{lemma}
\label{lemma-colim-and-lim-of-duals}
\begin{slogan}
Trivial duality for systems of perfect objects.
\end{slogan}
Let $A$ be a ring. Let $(K_n)_{n \in \mathbf{N}}$ be a system of
perfect objects of $D(A)$. Let $K = \text{hocolim} K_n$ be the derived colimit
(Derived Categories, Definition \ref{derived-definition-derived-colimit}).
Then for any object $E$ of $D(A)$ we have
$$
R\Hom_A(K, E) = R\lim E \otimes^\mathbf{L}_A K_n^\vee
$$
where $(K_n^\vee)$ is the inverse system of dual perfect complexes.
\end{lemma}

\begin{proof}
By Lemma \ref{lemma-dual-perfect-complex} we have
$R\lim E \otimes^\mathbf{L}_A K_n^\vee =
R\lim R\Hom_A(K_n, E)$
which fits into the distinguished triangle
$$
R\lim R\Hom_A(K_n, E) \to
\prod R\Hom_A(K_n, E) \to
\prod R\Hom_A(K_n, E)
$$
Because $K$ similarly fits into the distinguished triangle
$\bigoplus K_n \to \bigoplus K_n \to K$ it suffices to show that
$\prod R\Hom_A(K_n, E) = R\Hom_A(\bigoplus K_n, E)$.
This is a formal consequence of (\ref{equation-internal-hom})
and the fact that derived tensor product commutes with direct sums.
\end{proof}

\end{document}
